%==============================================================================
%  CHAPTER 2 — BUSINESS UNDERSTANDING AND TECHNICAL STUDY
%==============================================================================
\chapter{Business Understanding and Technical Study}
\label{chap:study}

\section{Introduction}
Before designing the platform, this chapter clarifies \emph{what} it must
achieve and \emph{on what} it should be built. It states the business
objectives and the system objectives that answer them, analyzes the existing
manual workflow against the \gls{mlops} maturity scale, studies the platforms
that already address this space, and closes with the benchmark of technologies
selected to build the solution.

% ---------------------------------------------------------------------------
\section{Business Objectives}
\label{sec:bo}
Six business objectives were fixed with the host company:
\begin{enumerate}[leftmargin=1.6em]
  \item \textbf{Train without owning infrastructure}: a data scientist should
        use cluster compute without any cloud or Kubernetes knowledge;
  \item \textbf{Never lose an experiment}: every run must be recorded and
        comparable;
  \item \textbf{Choose the best model objectively}: model selection must rest
        on evaluation metrics, not memory or habit;
  \item \textbf{Ship a model like software}: versioned, promoted, deployed in
        a repeatable way;
  \item \textbf{Let external applications consume models}: a deployed model
        must be usable by any website or service;
  \item \textbf{Operate and trust the system}: usage, health, and failures
        must be visible and explainable.
\end{enumerate}

% ---------------------------------------------------------------------------
\section{System Objectives}
\label{sec:so}
Each business objective maps to a delivered system capability, summarized in
Figure~\ref{fig:bodso}.

\begin{figure}[H]\centering
  \includegraphics[width=\textwidth,keepaspectratio]{bodso}
  \caption{From business objectives to delivered system capabilities}
  \label{fig:bodso}
\end{figure}

Concretely: any \texttt{.py} / \texttt{.ipynb} / \texttt{.zip} runs in an
isolated, resource-limited Kubernetes pod (1); MLflow autolog captures
parameters, metrics, and artifacts of every run with zero code changes (2); a
leaderboard groups runs by model family and ranks them on evaluation metrics
only (3); the registry versions and stages models, and one click deploys a
version as a KServe \texttt{InferenceService} (4); a public endpoint secured
by hashed \gls{api} keys serves external applications without a platform
account (5); and serving telemetry, model-evolution \glspl{kpi}, failure
diagnostics, and an administration area make the system operable (6).

% ---------------------------------------------------------------------------
\section{Analysis of the Existing Workflow}
\label{sec:existing}
Google's reference architecture for \gls{mlops}~\cite{google_mlops} classifies
organizations into three maturity levels, a useful yardstick for evaluating
both the current situation and the target:
\begin{itemize}[leftmargin=1.4em]
  \item \textbf{Level 0, manual process}: every step is manual and
        script-driven, with a complete disconnection between the data
        scientists who produce models and the engineers who serve them. The
        notebook-on-a-laptop workflow of Chapter~\ref{chap:context} sits
        squarely at this level.
  \item \textbf{Level 1, ML pipeline automation}: training is automated as a
        re-executable pipeline; experiments are tracked; models are versioned
        in a registry and deployed through a repeatable process.
  \item \textbf{Level 2, CI/CD pipeline automation}: the pipeline itself is
        built, tested, and deployed automatically; retraining can be triggered
        by new data or detected drift.
\end{itemize}
The platform built in this project takes a team from level~0 to level~1, and
part of the way to level~2: the application itself is continuously built and
deployed through \gls{cicd}, while retraining remains user-triggered
(Figure~\ref{fig:maturity-ladder}).

\begin{figure}[H]\centering
\begin{tikzpicture}[
  lvl/.style={rounded corners=6pt, align=center, text width=3.5cm,
    minimum height=2.4cm, inner sep=6pt, font=\small, softshadow}]
  % three maturity levels, left to right, same height
  \node[lvl, draw=diagslate!40, fill=diagslate!6, text=diagslate] (l0) at (0,0)
    {{\scriptsize\bfseries LEVEL 0}\\[3pt]{\bfseries Manual process}\\[4pt]
     {\footnotesize Notebooks, scripts,\\ hand-made deploys}};
  \node[lvl, draw=diaggreen!65, line width=1.2pt, fill=diaggreen!12, text=diagslate] (l1) at (5,0)
    {{\scriptsize\bfseries\textcolor{diaggreen!45!black}{LEVEL 1}}\\[3pt]{\bfseries ML pipeline automation}\\[4pt]
     {\footnotesize Tracked runs, registry,\\ repeatable deploys}};
  \node[lvl, draw=diagblue!40, dashed, fill=diagblue!4, text=diagslate] (l2) at (10,0)
    {{\scriptsize\bfseries\textcolor{diagblue}{LEVEL 2}}\\[3pt]{\bfseries CI/CD\\ automation}\\[4pt]
     {\footnotesize Automated retraining\\ on new data}};
  \draw[-{Latex[length=2.6mm]}, line width=1pt, draw=diagslate!55] (l0.east) -- (l1.west);
  \draw[-{Latex[length=2.6mm]}, line width=1pt, draw=diagslate!55] (l1.east) -- (l2.west);
  % what the platform delivers, annotated below, no rotated text
  \draw[-{Latex[length=3mm]}, line width=2pt, draw=diaggreen!70] (0,-1.85) -- (5,-1.85);
  \node[font=\scriptsize\bfseries, text=diaggreen!45!black, anchor=north] at (2.5,-2.05)
    {This platform: Level 0 $\rightarrow$ Level 1};
  \draw[-{Latex[length=3mm]}, line width=1.4pt, draw=diagblue!55, dashed] (5,-1.85) -- (10,-1.85);
  \node[font=\scriptsize, text=diagslate, anchor=north] at (7.5,-2.05)
    {partially (application via CI/CD)};
\end{tikzpicture}
\caption{MLOps maturity ladder and the step delivered by the platform}
\label{fig:maturity-ladder}
\end{figure}

% ---------------------------------------------------------------------------
\section{State of the Art}
\label{sec:platforms}
Several commercial and open-source systems already address this space, from
fully managed cloud services (Amazon SageMaker~\cite{sagemaker}, Google Vertex
AI~\cite{vertexai}, Azure Machine Learning~\cite{azure_ml}, Databricks) to a
commercial training platform (AIchor~\cite{aichor}) and an open-source toolkit
(Kubeflow).

The managed clouds look complete but demand a cloud expert and significant
setup (accounts, permissions, storage, SDKs) before training a single model;
they are also closed and locked to one provider. The closest system to this
project's positioning is \textbf{AIchor}, developed by InstaDeep (a company
founded in Tunisia): it distributes training workloads on a cluster through a
GitOps workflow, but its lifecycle \emph{stops at training}, it is a
proprietary product, and it still requires ``a data scientist with Docker
skills''.

The platform built here removes these barriers: it covers the \emph{full}
lifecycle through serving and public consumption by \gls{api} key, is
assembled entirely from open-source components, and needs nothing more than a
browser. Table~\ref{tab:capability-matrix} compares the platforms feature by
feature.

\begin{table}[H]\centering
\caption{Feature-by-feature capability comparison}
\label{tab:capability-matrix}
\renewcommand{\arraystretch}{1.45}
\footnotesize
\begin{tabularx}{\textwidth}{|>{\raggedright\arraybackslash}p{3.15cm}|*{4}{>{\centering\arraybackslash}X|}>{\columncolor{primary!12}\centering\arraybackslash}X|}
\hline
\rowcolor{primary!15}
\textbf{Capability} & \textbf{SageMaker $\cdot$ Vertex $\cdot$ Azure~ML} & \textbf{Data\-bricks} & \textbf{AIchor} & \textbf{Kube\-flow} & \textcolor{primary}{\textbf{This plat\-form}}\\
\hline
Full ML lifecycle (train $\to$ deploy $\to$ consume) & \capyes & \cappart & \capno & \cappart & \capyes\\ \hline
Automatic experiment tracking & \capyes & \capyes & \cappart & \cappart & \capyes\\ \hline
Model registry \& versioning & \capyes & \capyes & \capno & \cappart & \capyes\\ \hline
One-click deployment (serving) & \capyes & \cappart & \capno & \cappart & \capyes\\ \hline
Public \gls{api} for external apps & \cappart & \cappart & \capno & \capno & \capyes\\ \hline
100\% open-source $\cdot$ no lock-in & \capno & \capno & \capno & \capyes & \capyes\\ \hline
Browser-only $\cdot$ no Git/Docker & \capno & \cappart & \capno & \capno & \capyes\\ \hline
\end{tabularx}

\vspace{0.35em}
{\small \capyes~Full support \quad \cappart~Partial \quad \capno~Not available
\hfill \textbf{\textcolor{primary}{Only this platform covers every capability.}}}
\end{table}

% ---------------------------------------------------------------------------
\section{Technology Benchmark and Selection}
\label{sec:tech-study}
For each concern of the platform, an open-source technology was selected; each
choice, its justification, and its logo follow.

% one centered logo with a numbered "Figure x.y: Name" caption below it
\newcommand{\techpic}[2]{\par\vspace{2pt}{\centering #1\par\vspace{1pt}%
  \captionof{figure}{#2}\par}\vspace{6pt}}
% two logos side by side, each with its own numbered caption
\newcommand{\techpair}[4]{\par\vspace{2pt}\noindent
  \begin{minipage}[t]{0.48\linewidth}\centering #1\par\vspace{1pt}\captionof{figure}{#2}\end{minipage}\hfill
  \begin{minipage}[t]{0.48\linewidth}\centering #3\par\vspace{1pt}\captionof{figure}{#4}\end{minipage}\par\vspace{6pt}}

\paragraph{Containerization and orchestration: Docker and Kubernetes.}
Containers~\cite{docker_docs} package an application with its exact
dependencies, and \gls{k8s}~\cite{kubernetes_docs} orchestrates them at scale.
Kubernetes is not only the deployment target but a core \emph{feature} of the
product: every training job and every deployed model runs in its own
resource-limited pod, which gives the platform isolation and elasticity. The
platform runs on \gls{aks} in production and on KinD locally, with no code
changes.
\techpic{
\includegraphics[height=18mm,keepaspectratio]{kubernetes}}{Kubernetes}

\paragraph{Pipeline orchestration: Kubeflow Pipelines and Argo Workflows.}
\gls{kfp}~\cite{kubeflow_pipelines} expresses a training workflow as a small
Python graph compiled into Argo Workflows~\cite{argo_workflows}, which runs
each step as an isolated pod with scheduling, retries, and log capture handled
by Argo. The backend submits one pipeline per run and follows its progress
through the \gls{kfp} \gls{api}, which is also how live logs reach the
browser. Apache Airflow was considered but targets scheduled data workflows
rather than on-demand, per-user \gls{ml} jobs on Kubernetes.
\techpair{
\includegraphics[height=17mm,keepaspectratio]{kubeflow}}{Kubeflow Pipelines}%
  {\includegraphics[height=17mm,keepaspectratio]{argo}}{Argo Workflows}

\paragraph{Experiment tracking and registry: MLflow.}
\textbf{MLflow}~\cite{mlflow_docs} records parameters, metrics, and artifacts
of a run through a single \texttt{mlflow.autolog()} call, and provides a
built-in \textbf{Model Registry} with Staging / Production / Archived stages.
The platform provisions one MLflow experiment per project and mirrors registry
state in its own database so the interface stays responsive.
\techpic{
\includegraphics[height=17mm,keepaspectratio]{mlflow}}{MLflow}

\paragraph{Model serving: KServe.}
\textbf{KServe}~\cite{kserve_docs} turns a declarative
\texttt{InferenceService} resource into a running Deployment, Service, and
Ingress backed by a framework-specific model server, used here in
\textbf{RawDeployment} mode. Deploying a model thus reduces to creating a
single Kubernetes resource pointing at the model artifact in object storage.
Seldon Core offers similar capabilities but with a heavier operational
footprint for a single-team platform.
\techpic{
\includegraphics[height=19mm,keepaspectratio]{kserve}}{KServe}

\paragraph{Object storage: MinIO.}
\textbf{MinIO}~\cite{minio_docs} is an \gls{s3}-compatible object store
deployed inside the cluster. It holds uploaded code, datasets, and all MLflow
artifacts, and is a drop-in for the \gls{s3} \gls{api} that MLflow, KServe,
and the AWS SDK expect.
\techpic{
\includegraphics[height=14mm,keepaspectratio]{minio}}{MinIO}

\paragraph{Database: PostgreSQL.}
A single \textbf{PostgreSQL}~\cite{postgresql_docs} instance stores both the
platform's own tables and MLflow's tracking tables, keeping operations simple.
The backend accesses it asynchronously through SQLAlchemy~2, so database I/O
never blocks the event loop.
\techpic{
\includegraphics[height=16mm,keepaspectratio]{postgresql}}{PostgreSQL}

\paragraph{Backend: FastAPI.}
\textbf{FastAPI}~\cite{fastapi_docs} is an asynchronous Python web framework
that derives request validation and an OpenAPI schema from type hints. Python
is the natural choice because the whole machine-learning ecosystem (MLflow,
the Kubeflow SDK, the Kubernetes client) is Python. Compared with Flask or
Django, FastAPI brings native async support and generated documentation with
no extra code.
\techpic{
\includegraphics[height=16mm,keepaspectratio]{fastapi}}{FastAPI}

\paragraph{Frontend: Angular.}
\textbf{Angular}~\cite{angular_docs} is a batteries-included \gls{spa}
framework; the project uses Angular~19 with standalone components and signals,
served by NGINX. Route guards enforce authentication and the administrator
role on the client side, mirroring the backend's checks.
\techpic{
\includegraphics[height=16mm,keepaspectratio]{angular}}{Angular}

\paragraph{Secure public exposure: Cloudflare Tunnel.}
The platform is published through a \textbf{Cloudflare
Tunnel}~\cite{cloudflare_tunnel}: an in-cluster agent opens an outbound-only
connection, so no public IP address or inbound port is exposed. Public traffic
is terminated at Cloudflare's edge over \gls{https} and passes through its
protections, including \gls{ddos} mitigation.
\techpic{
\includegraphics[height=15mm,keepaspectratio]{cloudflare}}{Cloudflare Tunnel}

\paragraph{CI/CD: GitHub Actions and Azure Container Registry.}
Every push to the main branch triggers a \textbf{GitHub Actions} workflow that
builds the backend and frontend images, pushes them to \textbf{\gls{acr}},
and upgrades the cluster deployment.
\techpair{
\includegraphics[height=17mm,keepaspectratio]{github_actions}}{GitHub Actions}%
  {\includegraphics[height=17mm,keepaspectratio]{acr}}{Azure Container Registry}

\medskip
Table~\ref{tab:software} summarizes the final technical choices, together with
the working environment.

\begin{table}[H]\centering
\renewcommand{\arraystretch}{1.2}
\caption{Final technical choices}\label{tab:software}
\small
\begin{tabularx}{\textwidth}{|>{\raggedright\arraybackslash}p{4.2cm}|X|}
\hline
\rowcolor{primary!15}
\textbf{Concern} & \textbf{Selected technology}\\
\hline
Backend             & FastAPI (Python~3.11), SQLAlchemy~2 (async), Pydantic~v2\\
\hline
Frontend            & Angular~19 (standalone components, signals), Tailwind\\
\hline
Database            & PostgreSQL~16 (platform + MLflow)\\
\hline
Object storage      & MinIO (\gls{s3}-compatible)\\
\hline
Experiment tracking & MLflow (tracking + registry)\\
\hline
Orchestration       & Kubernetes (\gls{aks} / KinD), \gls{kfp}, Argo Workflows\\
\hline
Model serving       & KServe (RawDeployment)\\
\hline
Networking          & Cloudflare Tunnel (\gls{tls} at the edge)\\
\hline
CI/CD               & GitHub Actions, \gls{acr}\\
\hline
Development         & Visual Studio Code, Git/GitHub, Docker\\
\hline
Modeling \& report  & \gls{uml} (PlantUML, draw.io), \LaTeX\\
\hline
Hardware            & Intel Core i7 / 16\,GB workstation; \gls{aks} cluster\\
\hline
\end{tabularx}
\end{table}

% ---------------------------------------------------------------------------
\section{Conclusion}
This chapter fixed the business and system objectives, situated the existing
manual workflow at maturity level~0, showed that no existing platform covers
every needed capability, and selected an open-source technology for each
concern. The next chapter designs the system built on these choices.
